\section{Variable dependencies on $I(t)$}
\label{sec:dependencies}

It happens to be possible to write all of the formulas $J(t)$,
$\text{Var}(\xt)$, $V(t)$ and $\text{Var}(\wt)$ as functions of the
non-rupture probability $I(t)$ alone. Time enters these expressions only
through $I$, so the whole family of curves is parametrised by a single
quantity moving monotonically from $1$ down to $b$; this is the viewpoint
adopted in Chapter~4b, where the single-cell formulas become the kernels of a
population model.

\subsection{$J(I)$ and $K(I)$}

Since $J = -\delta^{-1}\,{\dd I}/{\dd t}$ (\ref{ddIJ}), and
${\dd I}/{\dd t} = \lambda (I - a)(I-b)$ (\ref{ddIIfact}), we can write $J$
in terms of $I$:
\begin{equation}
\label{JI}
J(I) = -\frac{\lambda}{\delta} (I - a)(I-b).
\end{equation}
Combining this with Theorem~\ref{thm:K}, the second moment likewise depends on
$I$ alone:
\begin{equation}
\label{KI}
K(I) = \frac{2\lambda^2}{\delta^2}(I-\kappa)(I - a)(I-b),
\end{equation}
with $\kappa = 1+\delta/2\lambda$.

\subsection{Mean and variance of $\xt$}

The variance of $\xt$ follows from $\text{Var}(\xt) = K(I) - J(I)^2$ written
entirely in terms of $I$:
\begin{align*}
\text{Var}(\xt) &= K(I) - {J(I)}^2\\
&=\frac{2\lambda^2}{\delta^2}\lrb{I-\kappa}\lrb{I-a}\lrb{I-b}
- \lrb{\frac{\lambda}{\delta}(I - a)(I-b)}^2,
\end{align*}
which, like $J$ and $K$, vanishes at either absorbing limit $I\to a$ or
$I\to b$ and is positive between them.

\subsection{Mean and variance of $\wt$}

As we have defined $V(t)$ in terms of $I(t)$ and $J(t)$, it is likewise
possible to use (\ref{JI}) to write $V$ as a function only of the non-rupture
probability:
\begin{equation}
\label{VI}
V(I) = \frac{\lambda - \mu}{\delta}\lrb{1-I} - J(I) + 1
= \frac{\lambda - \mu}{\delta}\lrb{1-I} + \frac{\lambda}{\delta} (I - a)(I-b) + 1.
\end{equation}
And given (\ref{VarY}), we can do the same for the variance of $\wt$:
\begin{equation}
\label{VarWI}
\text{Var}(\wt) = \frac{2(\lambda-\mu)}{\delta} V(I) - \frac{\lambda+\mu}{\delta}I - K(I) + \frac{\lambda+\mu}{\delta} + 1 - {V(I)}^2.
\end{equation}
Written out in full, with (\ref{JI}), (\ref{KI}) and (\ref{VI}) substituted:
\begin{multline*}
\text{Var}(\wt) = \frac{2(\lambda - \mu)}{\delta} \lrbs{ \frac{\lambda - \mu}{\delta}\lrb{1-I} + \frac{\lambda}{\delta} (I - a)(I-b) + 1} \\
- \frac{\lambda+\mu}{\delta} I
-\frac{2\lambda^2}{\delta^2}\lrb{I-\kappa}\lrb{I-a}\lrb{I-b}
+ \frac{\lambda + \mu}{\delta} + 1 \\ - \lrbs{\frac{\lambda - \mu}{\delta}\lrb{1-I} + \frac{\lambda}{\delta} (I - a)(I-b) + 1}^2.
\end{multline*}
The late-time value of \eqref{VarWI} is the variance of the burst size (over
all realisations, dead ones contributing zero); at $I=b$ with $\mu>0$ this
remains finite, and when $\mu=0$ it reduces to the variance of the geometric
law of Section~\ref{sec:variance}, namely $(1-s)/s^2$ with
$s=\delta/(\lambda+\delta)$.
